%% kapa: exam=1 round=25 year=2014 period=2 form=A subject=appraisal_law subject_ko=부동산관계법규 seq=11 number=11
%% source: Q-Net 감정평가사 기출문제 (cst00309) artlSeq=5101207 file=extracted/1008/5101207/2014년 제25회 감정평가사 1차시험 기출문제/2교시(A형).hwp
%% conversion: kordoc 4.24.51 -> tools/gen1.py (auto); review=proofread-sonnet; flags=none
\begin{question}{11}
국토의 계획 및 이용에 관한 법령상 도시ㆍ군계획시설사업의 시행에 관한 설명으로 옳지 않은 것은?
\begin{choices}
  \item 도시ㆍ군계획시설결정의 고시일로부터 2년 이내에 사업의 시행에 대한 단계별 집행계획이 수립되어야 한다.
  \item 시행자는 사업을 효율적으로 추진하기 위하여 필요하다고 인정되면 사업시행대상지역을 둘 이상으로 분할하여 사업을 시행할 수 있다.
  \item 사업으로 인하여 기반시설의 설치가 필요한 경우 사업의 시행자인 지방자치단체는 그 이행의 담보를 위한 이행보증금을 예치하여야 한다.
  \item 시행자는 사업시행을 위하여 특히 필요하다고 인정되면 도시ㆍ군계획시설에 인접한 토지를 일시 사용할 수 있다.
  \item 사업의 착수예정일의 변경을 내용으로 하는 실시계획 변경인가를 하는 경우에는 그에 대한 공고 및 열람을 하지 아니할 수 있다.
\end{choices}
\end{question}
\answer{3}
